\documentclass[10pt,a4paper]{article}
\usepackage[margin=0.8in]{geometry}
\usepackage{amsmath,amssymb,hyperref}
\hypersetup{colorlinks=true,urlcolor=blue,linkcolor=blue}
\setlength{\parindent}{0pt}\setlength{\parskip}{5pt}
\newcommand{\Z}{\mathbb Z}
\newcommand{\HC}{\operatorname{HC}}
\begin{document}
\begin{center}
{\Large Integral cyclic cohomology is not two-periodic from degree zero}\\[4pt]
TLMC 00000001537 \quad lizaixi01 \quad 24 October 2026
\end{center}

\textbf{The printed assertion and the convention.}
The original source defines cyclic cohomology $\HC^n(R)$ and asserts period two
for $R=\Z[X]/(X^2)$. We use ordinary cyclic cohomology over the coefficient
ring $\Z$, indexed from $n=0$. In this standard reading, period two would
in particular require an isomorphism of abelian groups
\[
  \HC^0(R)\cong\HC^2(R).
\]
We disprove this necessary consequence: degree zero has no nonzero element
killed by addition to itself, whereas degree two has such an element.
No assertion concerns periodic cyclic cohomology $\operatorname{HP}$, a change
to rational coefficients, or an unspecified eventual periodicity.

\textbf{The standard cohomology model.}
Loday's Definition 2.4.1 defines ordinary cyclic cohomology from the dual
cyclic bicomplex, with
$C^n(A)=\operatorname{Hom}_{\Z}(A^{\otimes(n+1)},\Z)$.
Section 2.4.3 identifies its cohomology with that of the normalized
$(b^*,B^*)$ total cochain complex. This is the standard model formalized here.
Its normalized cochains are the full multilinear functionals that vanish
whenever some argument $a_i$, $i\ge1$, equals the unit.
The cyclic-invariant single complex in Section 2.4.2 requires rational
coefficients for the stated comparison and is not used.

In Lean, $R$ is the native \texttt{AdjoinRoot} of $X^2\in\Z[X]$.
The universal properties construct an actual algebra equivalence
\[
  R\simeq_{\Z}D=\Z[\varepsilon],\qquad \varepsilon^2=0,
\]
where $D$ is the native dual-number algebra. Its full normalized multilinear
spaces are used in every degree. The native tensor universal property
identifies full cochains with linear maps from the genuine tensor powers.
The span of tensors with a positive-position unit is an actual submodule;
the dual of its native quotient is proved equivalent to the normalized
cochain space. Also $D/(\Z\cdot1)\simeq_{\Z}\Z$, through the actual
epsilon-coordinate map.

\textbf{All normalized cochains, not selected test functionals.}
Write $z=c(z)+e(z)\varepsilon$. For a normalized $f\in\overline C^n(D)$,
set
\[
 a=f(1,\varepsilon,\ldots,\varepsilon),\qquad
 b=f(\varepsilon,\ldots,\varepsilon).
\]
Multilinearity and normalization give, on every input,
\[
 f(z_0,\ldots,z_n)=
 \bigl(a\,c(z_0)+b\,e(z_0)\bigr)\prod_{i=1}^n e(z_i).
 \tag{1}
\]
Conversely this formula gives a normalized cochain for every $(a,b)\in\Z^2$.
The proof classifies the full space by the genuine two-element module basis,
so $(a,b)$ are actual evaluation coordinates, not assigned cohomology groups.

\newpage
\textbf{The genuine operators in every degree.}
The Hochschild faces multiply adjacent arguments; the final face multiplies
the last and first arguments in that order. With their alternating signs,
\[
 (b^*f)(z_0,\ldots,z_{n+1})=
 \sum_{j=0}^n(-1)^j f(z_0,\ldots,z_jz_{j+1},\ldots,z_{n+1})
 +(-1)^{n+1}f(z_{n+1}z_0,z_1,\ldots,z_n).
 \tag{2}
\]
Each actual face is bundled as a multilinear map using the products in
the fibers of its slot map. For an input with a positive-position unit,
all terms vanish except the two equal neighboring faces with opposite signs.
At the last position, the final face is the cancelling neighbor. Thus
(2) truly preserves normalization on all inputs.

For Connes' map use the full expression $(1-t)sN$: $s$ inserts the unit,
$t$ is the signed right cyclic rotation, and $N=1+t+\cdots+t^n$.
The finite cyclic sum is checked against genuine iterates of the rotation.
After evaluation by a normalized cochain, the $tsN$ terms vanish because
their second argument is the unit. Substituting (1) into these actual
operators gives the following full coordinate formulas:
\[
 b^*:\overline C^n\longrightarrow\overline C^{n+1},\quad
       (a,b)\longmapsto\bigl((1+(-1)^{n+1})b,0\bigr),
 \tag{3}
\]
\[
 B^*:\overline C^{n+1}\longrightarrow\overline C^n,\quad
       (a,b)\longmapsto(0,c_n a),\qquad
 c_n=\sum_{i=0}^n(-1)^{ni}=
 \begin{cases}n+1&n\text{ even},\\0&n\text{ odd}.\end{cases}
 \tag{4}
\]
The Lean proof verifies the standard formulas on every input; (3)--(4)
are derived from those formulas.

\textbf{The full native total complex and the actual obstruction.}
In degree $n$ take the finite product
\[
 T^n=\prod_{p=0}^{\lfloor n/2\rfloor}\overline C^{n-2p}(D).
\]
The total differential sums each genuine $b^*$ component in its column
and each genuine $B^*$ component into the next column. Equations (3)--(4)
give $d^{2k}=0$ in all degrees, hence $d^{n+1}d^n=0$.
This defines a native \texttt{CochainComplex} in \texttt{ModuleCat} $\Z$;
cohomology is its native homology object.

There is no incoming differential in degree zero, and $d^0=0$.
The native homology API therefore identifies $\HC^0(D)$ with $T^0\simeq\Z^2$.
In degrees one and two, the same actual differential is
\[
 d^1(a,b)=\bigl((2b,0),(0,a)\bigr),\qquad d^2=0.
 \tag{5}
\]
Consequently native $\HC^2(D)$ is identified with
$T^2/\operatorname{im}d^1$. Let $u=((1,0),(0,0))\in T^2$.
Its class is nonzero: the first coordinate of every boundary is an even
integer. But $2u=d^1(0,1)$, so $[u]+[u]=0$.
These are actual cycle, boundary and quotient-class calculations.
Since $\Z^2$ has no nonzero element with that property, no abelian-group
isomorphism $\HC^0(D)\simeq\HC^2(D)$ exists.

\newpage
\textbf{Back to the exact source algebra.}
The proof does not stop with an algebra equivalence appended to a calculation
on another object. For the actual native $R=\Z[X]/(X^2)$ it constructs all
normalized multilinear cochain spaces and their tensor--Hom maps.
Precomposition with the algebra equivalence and its inverse gives actual
linear equivalences in every degree. The Hochschild and Connes operators
on $R$ satisfy the standard formulas on every $R$-valued input; preservation
of multiplication and the unit proves their naturality.
The full total complexes are isomorphic, and the native homology API supplies
the resulting cohomology isomorphisms.

Thus the final theorem is directly about the native cohomology of the full
normalized total complex on $R$:
\[
 \begin{gathered}
 \forall x\in\HC^0(R),\quad x+x=0\Longrightarrow x=0,\\
 \exists y\in\HC^2(R),\quad y\ne0\ \text{and}\ y+y=0,\\
 \text{there is no additive-group equivalence }\HC^0(R)\simeq\HC^2(R).
 \end{gathered}
\]
This refutes the printed period-two assertion under the stated ordinary
integral, degree-zero convention. A classification of all coefficient
groups is unnecessary for this obstruction and is not asserted.

\textbf{What is formalized and what identifies the model.}
The general canonical comparison between the dual cyclic bicomplex and its
normalized $(b^*,B^*)$ model is the published standard identification in
Loday, Section 2.4.3. Lean formalizes the latter model explicitly: the full
cochain spaces and normalization, the genuine operators and their input
formulas, all diagonal total spaces and differentials, the native complex,
its actual homology, and the source-algebra transport. It does not assume
the proposed coefficient groups or infer cohomology by dualizing a list of
homology groups. Rudman's paper is used only for the standard chain-level
operator formulas, not as a cohomology computation.

\textbf{The supplied Lean project.}
Key declarations, all in namespace \texttt{TLMC1537}, include:

\begin{tabular}{@{}ll@{}}
\texttt{quotientDualEquiv} & \texttt{normalizedTensorDualEquiv}\\
\texttt{hochschild\_full\_formula} & \texttt{connes\_full\_formula}\\
\texttt{cyclicCochainComplex} & \texttt{HC0\_no\_add\_two\_torsion}\\
\texttt{torsionHC2\_nonzero} & \texttt{sourceHochschild\_full\_formula}\\
\texttt{sourceConnes\_full\_formula} & \texttt{sourceCyclicCochainComplex}\\
\texttt{complete\_counterexample} &
\end{tabular}

The complete original bilingual source is preserved separately.
The project pins Lean 4.33.0 and Mathlib commit
\texttt{db584cd6d46c92f209a44c0f1c829460d327499d}.
Command outputs, exhaustive namespace axiom checks, source hashes and the
official fresh-kernel receipt accompany the submission. Only the standard
axioms \texttt{propext}, \texttt{Classical.choice}, \texttt{Quot.sound}
are permitted. Production resource observations are diagnostic and do not
establish a controlled research-efficiency gain. Organizer review is pending.

\textbf{Primary references.}

J.-L. Loday, \emph{Cyclic Homology}, second edition, Grundlehren der
mathematischen Wissenschaften 301, Springer, 1998, printed pp.~72--73,
Definition 2.4.1 and Section 2.4.3.
\href{https://doi.org/10.1007/978-3-662-11389-9}{DOI: 10.1007/978-3-662-11389-9}.

E. Rudman, \emph{The cyclic homology of
$k[x_1,x_2,\ldots,x_d]/(x_1,x_2,\ldots,x_d)^2$},
arXiv:2010.08845v1, 2020, printed pp.~2--3, Definitions 2.1--2.5.
\href{https://arxiv.org/abs/2010.08845v1}{arXiv:2010.08845v1}.

\end{document}
